\Askhsh{34773}{2}
\begin{ylh}{1}
    \item 3.1. Είδη και στοιχεία τριγώνων
    \item 3.2. 1ο Κριτήριο ισότητας τριγώνων
    \item 3.4. 3ο Κριτήριο ισότητας τριγώνων.
\end{ylh}
% ===================================================================
%  Αρχείο: 34773.tex (Fragment)
%  Αυτόματη μετατροπή μέσω doc2latex.py
% ===================================================================

ΘΕΜΑ 2

Δίνεται ισοσκελές τρίγωνο ΑΒΓ (ΑΒ = ΑΓ) και Κ εσωτερικό σημείο του τριγώνου τέτοιο ώστε ΚΒ = ΚΓ.

\begin{alist}
\item Να αποδείξετε ότι:
\begin{rlist}
\item τα τρίγωνα ΒΑΚ και ΚΑΓ είναι ίσα, \monades{12}
\item η ΑΚ είναι διχοτόμος της γωνίας Β$\widehat{Α}$Γ. \monades{6}
\end{rlist}

\item Αν η προέκταση της ΑΚ τέμνει την ΒΓ στο Ε, τότε να δείξετε ότι η ΚΕ είναι διάμεσος του τριγώνου ΒΚΓ. \monades{7}

\begin{center}
\begin{tikzpicture}[scale=1]
\tkzDefPoint(0,0){B}
\tkzDefPoint(4,0){C}
\tkzDefTriangle[two angles = 70 and 70](B,C) \tkzGetPoint{A}
\tkzDefMidPoint(B,C) \tkzGetPoint{E}
\tkzDefPointBy[homothety=center E ratio 0.4](A) \tkzGetPoint{K}

\draw[l3] (A)--(B)--(C)--cycle;
\draw[l3] (A)--(E);
\draw[l3] (B)--(K)--(C);

\tkzDrawPoints(A,B,C,E,K)
\tkzLabelPoint[above](A){$A$}
\tkzLabelPoint[below](B){$B$}
\tkzLabelPoint[below](C){$\Gamma$}
\tkzLabelPoint[below](E){$E$}
\tkzLabelPoint[left](K){$K$}

\tkzMarkSegments[mark=s|](A,B A,C)
\tkzMarkSegments[mark=s||](K,B K,C)
\end{tikzpicture}
\end{center}
\end{alist}
